\scanpage{62}

\mspage{3}

Remplacant $A$ par $A'$, il reste à prouver que si $m/_{m^2}\otimes K\longrightarrow n/_{n^2}$ est surjectif \inserted{i é $K=B\otimes k$} on peut trouver $C$ tel que $K=C\otimes k$\inserted{ d'où $\operatorname{gr}C=K\otimes_k\operatorname{gr}A$}. On va manger $K$ par étapes

a) soit $(x_i)_{i\in I}$ une base de transcendance de $K$ sur $k$ ; \deleted{et} on relève les $x_i$ en des éléments de $B$, définissant ainsi $A(x_i)\longrightarrow B$. $A(x_i)$ est la limite inductive des localisées, au point générique de la fibre spéciale, des $A[X_i]_{i\in J}$, $J$ fini ds $I$. (cf Nagata, local rings pg 17 et 18). On peut soit utiliser $0_{\overline{\underline{\mathrm{III}}}}$ 10.3.1.3 pour vérifier que cet anneau est noethérien, soit le compléter et utiliser le critère de platitude pour vérifier que son \unsure{complété|complétée} est plat sur $A$

b) on s'est ainsi réduit à supposer $K$ algébrique sur $k$ et, si on y tient, à supposer $A$ complet. On trouve sans histoire $A'$ plat sur $A$, non ramifié, de corps résiduel la cloture séparable de $k$ dans $K$, et $A'\longrightarrow B$. Reste à traiter le cas où $K/k$ est purement inséparable, et on va utiliser une propriété ``d'intersection complète'' des corps, que je n'ai pas su bien dégager

c) on suppose : $K/k$ purement inséparable, $A$ et $B$ complets, $\deleted{\unsure{B}}K=\deleted{\unsure{A}}B\otimes k$.

soit $K_n$ le sous corps de $K$ formé des $x$ tels que $x^{p^n}\in k$ ($K_n=K\cap k^{p^{-n}}$), et soit $x_n^\alpha$ une $p$-base de $K_n$ sur $K_{n-1}$ ($n\geq1$). $K$ est \deleted{\unsure{le}} l'anneau engendré par des éléments $x_n^\alpha$ soumis à des relations $(x_n^\alpha)^p=P_n^\alpha(x_m^{\deleted{\unsure{\alpha}}\beta},m<n)$. \inserted{($P$ à coefficients dans $k$)}

Relevons les $x_n^\alpha$ dans $B$. On aura des relations

\[
(x_n^\alpha)^p=P_n^\alpha(x_m^\beta,m<n)+q_n^\alpha\quad\text{avec }q_n^\alpha\text{ dans l'idéal maximal }(P\text{ à \uline{coeff ds $A$}})
\]

par hypothèse, $m/_{m^2}\otimes_k K\longrightarrow n/_{n^2}$ est surjectif - D'où

\[
(x_n^\alpha)^p=P_n^\alpha(x_m^\beta,m<n)+Q_{n1}^\alpha(x_m^\beta)+{\unsure{r|z}}_n^\alpha
\quad\text{avec }P\ (\text{resp }Q)\text{ à coefficients dans}
\]

$A$ (resp $m$) et par récurrence et passage à la limite

\margin{$Q_{ni}^\alpha$ \unsure{linaires|linéaires|linéaire} en les $X_m^\beta$ \ldots\\ $\unsure{\bigcirc\ i|\text{Où}|(\ )^i}$ ?}

\[
\boxed{(x_n^\alpha)^p=P_n^\alpha(x_m^\beta,m<n)+\sum_{1}^{\infty}Q_{ni}^\alpha(x_m^\beta)}
\]

où $P_n^\alpha$ (resp $Q_{ni}^\alpha$) est un polynome en les $x_m^\beta$ à coefficients dans $A$ (resp dans $m^{\text{\textcircled{\scriptsize$n$}}}$)

Pour $C$, il suffira maintenant de prendre l'anneau complet engendré par des générateurs $x_n^\alpha$ soumis aux relations précédentes. Pour une définition rigoureuse, disons qu'on prend la limite projective des $A/m^{\deleted{\unsure{n}}k}$-algèbres engendrées par des générateurs $x_n^\alpha$ soumis aux relations

\[
(x_n^\alpha)^p=P_n^\alpha(x_m^\beta,m<n)+\sum_{1}^{\deleted{\unsure{n}}k-1}Q_{n\unsure{l|i}}^\alpha(x_m^\beta). \tag*{(1)}
\]

Il reste à voir qu'on obtient ainsi une $A/m^k$-algèbre plate.

Pour cela, il suffit vérifier que le quotient de $A[X_n^\alpha]$ par un nombre fini de
% archive sheet number: 61
